\section{Task, Data and Protocol}
\label{sec:task}
\FigureSystem

\paragraph{Task.}
Each item is a triple $x = (q, Q, a)$ from a U.S.\ presidential interview \citep{thomas-etal-2024-never}: a sub-question $q$ (median 15 tokens), the journalist's full turn $Q$ it was split from (62 tokens) and the President's whole answer $a$ (266 tokens).
Subtask~2 (\Stwo{}) asks for one of nine evasion types $\mathcal{Y}_9$, describing how $a$ responds to \emph{that} sub-question, and each type fixes one of three clarity levels for Subtask~1 (\Sone{}, Table~\ref{tab:taxonomy}).
Like the top systems, we predict the nine-way type and map it to the clarity level.

\begin{table}[t]
\centering
\small
\setlength{\tabcolsep}{3.5pt}
\begin{tabular}{@{}llrr@{}}
\toprule
Evasion type (\Stwo{}) & Clarity level (\Sone{}) & Train \% & Dev sets \\
\midrule
Explicit            & Clear Reply     & 30.5 & 115 \\
\midrule
Implicit            & Ambivalent      & 14.2 & 99 \\
Dodging             & Ambivalent      & 20.5 & 96 \\
General             & Ambivalent      & 11.2 & 113 \\
Deflection          & Ambivalent      & 11.0 & 51 \\
Partial/half-answer & Ambivalent      & 2.3  & 14 \\
\midrule
Declining to answer & Clear Non-Reply & 4.2  & 17 \\
Claims ignorance    & Clear Non-Reply & 3.5  & 15 \\
Clarification       & Clear Non-Reply & 2.7  & 4 \\
\bottomrule
\end{tabular}
\caption{The taxonomy. Train \% is the share of the 3,448 training labels. Dev sets is the number of the 308 dev items whose reference set contains the type.}
\label{tab:taxonomy}
\end{table}

\paragraph{Data.}
Train has 3,448 items with one label each.
Dev has 308 items with three annotators: 125 are unanimous, 150 split two to one and 33 receive three labels, so the reference set $G_i$ of an item holds one to three types.
The 237 test items have two annotators, but their labels were never released, so every comparison is on dev, and we also rescore dev against each of its three two-annotator subsets (the test regime).
One property shapes the whole paper: 69\% of training rows share their answer with another sub-question of the same turn (\emph{siblings}).

\paragraph{Metric.}
\Stwo{} is scored by multi-reference macro-F1: a prediction is a true positive if it lies in $G_i$, and every class counts equally.
It rewards both the \emph{in-set rate} and \emph{coverage} of rare classes.
Always predicting Explicit scores 0.060, TF-IDF with logistic regression 0.257, and a human annotator against the other two 0.684.

\paragraph{Protocol.}
Epochs, calibration and routing thresholds are chosen on the \emph{train slice}, a fixed stratified 10\% of train (345 rows).
The few parameters that must be fitted on dev are chosen by 5-fold nested cross-validation, and we report the first split and the mean over ten splits.
A \emph{system} averages seeds and applies logit adjustment \citep{menon2021longtail}, $\hat y = \arg\max_c \bar p(c \mid x)/\pi_c^{\tau}$ with $\pi_c$ the train prior and $\tau$ the only fitted parameter.
Claims rest on ten paired seeds, and a change must first pass a three-seed screen (+0.015, at least two of three seeds up).
Before each experiment we registered its hypotheses, predictions and decision rules. Later changes were logged as amendments before results were seen.
From \S\ref{sec:measure} on, every label comes from one \emph{frozen} system, so no parameter changes between experiments.
